\chapter{Arquitectura: capas de cizalla}
\label{sec:paper-arq}

\section{Idea central}

\emph{Shearing
layers}~\cite{brand1994how} organiza un sistema por velocidades de
cambio. Lo que cambia rápido (pieles, superficies, integraciones)
no debe acoplar a lo que cambia lento (esqueletos, motores). En
software, la consecuencia práctica es dirigir todas las
dependencias \emph{desde} lo rápido \emph{hacia} lo lento, y
prohibir dependencias en sentido contrario. Una arista
núcleo$\to$plataforma, si existe, es una fuente de fricción
estructural: el núcleo se ve forzado a moverse a la velocidad de
la superficie.

GraphHunter adopta tres capas con esta disciplina:

\begin{enumerate}
    \item \textbf{Core} (``esqueleto''): tipos del grafo, operaciones
    estructurales, scoring, matcher FFI, clasificador GNN. Cambia
    por meses.
    \item \textbf{Platform} (``carne''): DSL, parsers, sources de
    ingesta, canonical OCSF, servidor HTTP, MCP TypeScript,
    servicios compartidos (catálogo, DLQ, drift monitor, mapping
    library). Cambia por semanas.
    \item \textbf{Apps} (``piel''): Tauri desktop, CLI. Cambia por
    días.
\end{enumerate}

\begin{figure}[H]
\centering
\begin{tikzpicture}[
    node distance=0.6cm and 0.8cm,
    layer/.style={draw, rounded corners, minimum width=11cm,
                  minimum height=1.2cm, align=center, font=\small\bfseries},
    apps/.style={layer, fill=red!10, draw=red!50!black},
    platform/.style={layer, fill=orange!10, draw=orange!60!black},
    core/.style={layer, fill=blue!10, draw=blue!50!black},
]
\node[apps]     (apps)      {Apps\\\footnotesize\mdseries (tauri · cli)};
\node[platform] (platform)  [below=of apps]
    {Platform\\\footnotesize\mdseries (api · dsl · parsers · sources · canonical · mcp · mapping-library)};
\node[core]     (core)      [below=of platform]
    {Core\\\footnotesize\mdseries (graph-engine · matcher-ffi · gnn · libgraphmatch C++)};
\draw[->, thick] (apps)      -- (platform);
\draw[->, thick] (platform)  -- (core);
\end{tikzpicture}
\caption{Las tres capas. Todas las aristas de dependencia apuntan
hacia abajo (piel$\to$carne$\to$esqueleto). No hay ciclos.}
\label{fig:paper-capas}
\end{figure}

\section{Crates y lenguajes}

El sistema está compuesto por 14 \emph{crates} Rust, un servidor
MCP en TypeScript, y una biblioteca C++ (\filepath{libgraphmatch})
con intrínsecas AVX2 para el matcher. La elección de lenguajes
sigue la misma lógica de shearing:

\begin{itemize}
    \item \textbf{C++} donde el costo de abstracción es inaceptable:
    intersección SIMD de posting-lists en el núcleo del matcher.
    \item \textbf{Rust} en todo el sistema: tipos fuertes, zero-cost
    abstractions, seguridad sin recolector. Es el tejido por
    defecto.
    \item \textbf{TypeScript} sólo para la superficie MCP: el
    protocolo tiene SDK oficial en TS y la audiencia es
    asistentes LLM con tooling JS.
\end{itemize}

\section{Contraste con SOLID uniforme}

Una crítica posible es: \emph{``ya existe SOLID, ¿para qué otro
modelo?''}. SOLID es un conjunto de reglas micro, aplicables a
clases individuales. Shearing layers es una regla macro: aplica a
\emph{módulos de alto nivel}. Son complementarios. La diferencia
relevante es la dirección de la presión de diseño: SOLID empuja a
abstraer \emph{por contrato}; shearing empuja a separar \emph{por
velocidad de cambio}. Un sistema bien SOLID puede seguir teniendo
dependencias núcleo$\to$superficie; shearing prohibe eso
explícitamente.

\section{Propiedades verificables}

La descomposición se validó con tres propiedades:

\begin{enumerate}
    \item \textbf{Cero ciclos}. Verificado con
    \code{cargo-modules}/\code{cargo-deps} y \code{cargo check}
    sobre el workspace. Un ciclo rompería la capacidad de aislar
    lo que cambia.
    \item \textbf{FFI aislado en una capa}. Sólo
    \code{matcher-ffi} en el core toca C++. El resto del código
    Rust no ve \code{extern "C"}, no maneja punteros crudos y no
    preocupa por ABI. La versión ABI se chequea en runtime
    (cap.~\ref{sec:paper-ffi}).
    \item \textbf{Canonical API \emph{sealed}}. Los tipos
    que cruzan la superficie HTTP/MCP se definen en una capa de
    DTOs; las mutaciones del grafo canonical no escapan.
\end{enumerate}

Las dos aristas residuales \code{core}$\to$\code{platform} (ver
cap.~\ref{sec:paper-limites}) no son ciclos pero invierten la
dirección deseada. Son deuda reconocida.

\section{Relación con los otros libros}

Este capítulo resume lo mínimo para situar las contribuciones
técnicas que siguen. La versión detallada del modelo, con el
inventario completo de crates y seis ADRs, vive en el libro
\emph{Arquitectura} de la documentación.
